%% ============================================================
%% 分片 07：原文第 3 节「Topological Field Theory From ${\EUN}=4$ Super Yang-Mills Theory」（原文 glshort.tex 起于第 1216 行）
%% 由 tools/harvmac2tex.py 从 arXiv 源码（hep-th/0604151v3）自动生成。
%% 只译英文散文；公式、编号、\label、\cite 一律不改。
%% ============================================================

\subsection{The Topological Lagrangian}

\label{sec:toplag}

Our goal now is to find a Lagrangian that possesses the
topological symmetry for any value of $t$ and reduces when $M$ is
flat to the Lagrangian of the underlying ${\EUN}=4$ super
Yang-Mills theory.  We could do this by hand, starting with \eqref{eq:ygo}\
for flat $M$ and asking what curvature dependent terms are needed
to maintain the topological symmetry when $M$ is not flat.  We
will follow a different approach.

A very useful first step is to compute the algebra generated on
the fields by the the supersymmetries $Q_\ell$ and $Q_\rr$, or
equivalently by $Q=uQ_\ell+vQ_\rr$.  For the fields $
A,\phi,\psi$, $\tilde\psi$, and $\sigma$, we compute
% 式 (3.35)
\setcounter{equation}{34}
\begin{equation}\label{eq:ibo}
\begin{aligned}\delta_T{}^2 A & = -i(u^2+v^2)(-D\sigma)\\
\delta_T{}^2\phi & = -i(u^2+v^2)[\sigma,\phi]\\ \delta_T{}^2 \psi
& = -i(u^2+v^2)[\sigma,\psi]\\ \delta_T{}^2\tilde\psi& =
-i(u^2+v^2)[\sigma,\tilde\psi]\\ \delta_T{}^2\sigma& = 0.\\\end{aligned}
\end{equation}

These results can be summarized by saying that if $\Phi$ is any
field of $\EUK\geq 0$, then
$\delta_T{}^2\Phi=-i(u^2+v^2){{\lie}}_\sigma(\Phi)$, where
${{\lie}}_\sigma(\Phi)$ is the first order change in $\Phi$ in a
gauge transformation generated by $\sigma$ (so $\lie$ stands for
Lie derivative). For example, ${{\lie}}_\sigma(A)=-D\sigma$, and
for the other fields here, ${{\lie}}_\sigma(\Phi)=[\sigma,\Phi]$.
The relation $\delta_T{}^2\Phi=-i(u^2+v^2){{\lie}}_\sigma(\Phi)$
can be expanded in powers of $u$ and $v$ and is equivalent to the
following:
% 式 (3.36)
\setcounter{equation}{35}
\begin{equation}\label{eq:uvu}
\begin{aligned}\delta_\ell^2\Phi=\delta_r^2\Phi&=-i{{\lie}}_\sigma(\Phi)\\
\{\delta_\ell,\delta_\rr\}\Phi& = 0.\\\end{aligned}
\end{equation}

If, however, one computes $\delta_T{}^2\Phi$ for a field of
$\EUK<0$, one does not get $-i(u^2+v^2){{\lie}}_\sigma(\Phi)$ in
an obvious way. In fact, the relation
$\delta_T{}^2\Phi=-i(u^2+v^2){{\lie}}_\sigma(\Phi)$ does hold for
all $\Phi$ in the twisted ${\EUN}=4$ super Yang-Mills theory, but
for some fields, it only holds upon using the equations of motion.

Construction of supersymmetric Lagrangians is generally much
easier if one can introduce ``auxiliary fields,'' which are extra
fields that can ultimately be eliminated by their equations of
motion if one so chooses, such that the algebra, in the present
case $\delta_T{}^2\Phi=-i(u^2+v^2){{\lie}}_\sigma(\Phi)$, is
satisfied without using equations of motion.   In theories with a
great deal of supersymmetry, such as ${\EUN}=4$ super Yang-Mills
theory, it can be very difficult to find suitable auxiliary
fields.   The present example, as first noted in \cite{ref53}, is a
case in which this problem arises.

However, the problem can be greatly alleviated by introducing an
auxiliary field $P$, a zero-form with values in the Lie algebra,
so as to close the algebra for $\eta$ and $\tilde \eta $ (but not
$\chi$). One takes the transformation laws of
$\bar\sigma,\eta,\tilde\eta$, and $P$ to be
% 式 (3.37)
\setcounter{equation}{36}
\begin{equation}\label{eq:yfog}
\begin{aligned}\delta_T\bar\sigma & = iu\eta +i v\tilde\eta
\\
                    \delta_T\eta & = vP+u[\bar\sigma,\sigma]\\
                    \delta_T\tilde\eta & =
                    -uP+v[\bar\sigma,\sigma]\\
                    \delta_T P & =
                    -iv[\sigma,\eta]+i u[\sigma,\tilde\eta].\\\end{aligned}
\end{equation}

These equations will reduce to those in section 第~\ref{sec:twisting}~节\ once we
impose the equations of motion -- which notably will set
$P=D^*\phi$.

Once we have a closed algebra on a set of fields -- in this case
all fields in the theory including $P$ but excluding $\chi$, since
we have not closed the algebra on $\chi$ -- it is generally a
simple matter to find the possible invariant Lagrangians.  In the
present case, we would like to find a partial action ${{\cmmib
I}}_0$ for the the fields on which we have closed the algebra
which obeys $\delta_\ell{{\cmmib I}}_0=\delta_\rr{{\cmmib
I}}_0=0$, and hence, for any $u,v$, obeys $\delta_T {{\cmmib
I}}_0=0$ where $\delta_T=u\delta_\ell+v\delta_\rr$. These
properties will hold if ${{\cmmib I}}_0=\delta_\ell\delta_\rr V$
for some gauge-invariant $V$. Indeed, if $V$ is gauge-invariant,
then ${{\lie}}_\sigma(V)=0$ and \eqref{eq:uvu}\ reduces to
$\delta_\ell^2V=\{\delta_\ell,\delta_\rr\}V=\delta_\rr^2V=0$.  So
$\delta_\ell\delta_\rr V$ will automatically be annihilated by
$\delta$.

In addition, we want to pick $V$ so that the equation of motion
for $P$ is $P=D^*\phi$, so as to be compatible with the formulas
of section 第~\ref{sec:twisting}~节.  So we take % 式 (3.38)
\setcounter{equation}{37}
\begin{equation}\label{eq:obo}
V={2\over e^2}\int_M
d^4x\sqrt g\left(-{1\over 2}\Tr\,\eta\tilde\eta -i\Tr\,\bar\sigma
D^*\phi\right),
\end{equation}
 and compute that ${\cmmib
I}_0=\delta_\ell\delta_\rr V$ is % 式 (3.39)
\setcounter{equation}{38}
\begin{equation}\label{eq:nobo}
\begin{aligned}{{\cmmib
I}}_0= {2\over e^2}\int_Md^4x\sqrt g\Tr& \left({1\over 2}P^2
-PD^*\phi+{1\over 2}[\bar\sigma,\sigma]^2-
D_\mu\bar\sigma\,D^\mu\sigma
-[\phi_\mu,\sigma][\phi^\mu,\bar\sigma]\right.\\ & + i\tilde\eta
D_\mu\tilde\psi^\mu+i\eta
D_\mu\psi^\mu-i\tilde\eta[\psi_\mu,\phi^\mu]+i\eta[\tilde\psi_\mu,\phi^\mu]
\\ &\left.-{i\over 2}[\sigma,\tilde\eta]\tilde\eta -{i\over
2}[\sigma,\eta]\eta+i[\bar\sigma,\psi_\mu]\psi^\mu
+i[\bar\sigma,\tilde\psi_\mu]\tilde\psi^\mu\right).\\\end{aligned}
\end{equation}
 The
Euler-Lagrange equation for $P$ is $P=D^*\phi$, as desired, and
one can make the replacement % 式 (3.40)
\setcounter{equation}{39}
\begin{equation}\label{eq:olgog}
\Tr\left(
P^2-2PD^*\phi\right)\to-\Tr (D^*\phi)^2,
\end{equation}
 if one wishes.

${{\cmmib I}}_0$ manifestly possesses the topological symmetry for
any value of $t=v/u$.  Moreover, as it is of the form
$\delta_\ell\delta_\rr V$, where the metric of $M$ does not enter
in the definition of $\delta_\ell$ and $\delta_\rr $ but only in
the choice of $V$, it also has the key property that leads to a
topological field theory: its dependence on the metric of $M$ is
of the form $\delta_\ell\delta_\rr (\dots)$.

The only reason that ${{\cmmib I}}_0$ is not a satisfactory action
is that it is degenerate; it does not contain the kinetic energy
for the gauge fields or any terms containing $\chi$. It does not
seem to be possible to complete the construction of the action
with a construction as simple as that above.  Though it is
possible to add auxiliary fields so as to close the algebra on
$\chi$, it does not seem to be possible to do this in a way that
incorporates both $\delta_\ell$ and $\delta_\rr$ and is useful for
understanding the appropriate twisted ${\EUN}=4$ action. (The fact
that a $t$-dependent topological invariant occurs on the right
hand side of \eqref{eq:ormob}\ appears to be an obstruction to this.)
Instead, we will fix a particular value of $t$ and consider only
the differential $\delta_t=\delta_\ell+t\delta_\rr$. We add an
auxiliary field $H$, which is to be a two-form with values in the
Lie algebra, and postulate % 式 (3.41)
\setcounter{equation}{40}
\begin{equation}\label{eq:hco}
\begin{aligned}\delta_t\chi& = H\\
                           \delta_t H & =
                           -i(1+t^2)[\sigma,\chi].\end{aligned}
\end{equation}

This is compatible with
% 式 (3.42)
\setcounter{equation}{41}
\begin{equation}\label{eq:kino}
\delta_t^2(\Phi)=-i(1+t^2){{\lie}}_\sigma(\Phi),
\end{equation}
 for
$\Phi=\chi,H$; this is the specialization of
$\delta_T^2=-i(u^2+v^2)\lie_\sigma$ for $(u,v)=(1,t)$. For \eqref{eq:hco}\
to agree with the transformation of $\chi$ found in section
第~\ref{sec:twisting}~节, the equations of motion will have to give
$H^+=\EUV^+(t)$, $H^-=t\EUV^-(t)$. We will construct the action to
ensure this.

As before, an action annihilated by $\delta_t$ can be ${\cmmib
I}_1=\delta_t V_1$, for any gauge-invariant $V_1$.  We pick
% 式 (3.43)
\setcounter{equation}{42}
\begin{equation}\label{eq:usefuld}
V_1={2\over e^2}\int_Md^4x \sqrt g{1\over
(1+t^2)}\Tr\left(\chi^+_{\mu\nu}\left({1\over
2}H^{+\mu\nu}-\EUV^+(t)^{\mu\nu}\right)
+\chi^-_{\mu\nu}\left({1\over
2}H^{-\mu\nu}-t\EUV^-(t)^{\mu\nu}\right)\right).
\end{equation}
 Then
% 式 (3.44)
\setcounter{equation}{43}
\begin{equation}\label{eq:usefuldud}
\begin{aligned}{{{\cmmib I}}}_1= {1\over e^2}\int_Md^4x
\sqrt g&\left({2\over (1+t^2)}\Tr\left({1\over
2}H^+_{\mu\nu}H^{+\mu\nu}-H^+_{\mu\nu}\EUV^+(t)^{\mu\nu}\right)\right.\\
& + {2\over (1+t^2)}\Tr\left({1\over
2}H^-_{\mu\nu}H^{-\mu\nu}-tH^-_{\mu\nu}\EUV^-(t)^{\mu\nu}\right)\\
&+
2\Tr\left(\chi^+_{\mu\nu}\left(iD\psi+i[\tilde\psi,\phi]\right)^{\mu\nu}
+\chi^-_{\mu\nu}\left(iD\tilde\psi-i[\psi,\phi]\right)^{\mu\nu}\right)\\
&+\Tr\left(i\chi^+_{\mu\nu}[\sigma,\chi^{+\mu\nu}]+
i\chi^-_{\mu\nu}[\sigma,\chi^{-\mu\nu}] \right)\biggr).\\\end{aligned}
\end{equation}
 Upon
integrating $H^{\pm}$ out of this, we can write the equivalent
action % 式 (3.45)
\setcounter{equation}{44}
\begin{equation}\label{eq:usefuldudd}
\begin{aligned}{{\cmmib I}}_1= {1\over
e^2}\int_Md^4x \sqrt g&\left(-{t^{-1}\over
(t+t^{-1})}\Tr\,\EUV^+(t)_{\mu\nu}\EUV^+(t)^{\mu\nu}-{t\over
(t+t^{-1})}\Tr\,\EUV^-(t)_{\mu\nu}\EUV^-(t)^{\mu\nu}\right)\\ &+
2\Tr\left(\chi^+_{\mu\nu}\left(iD\psi+i[\tilde\psi,\phi]\right)^{\mu\nu}
+\chi^-_{\mu\nu}\left(iD\tilde\psi-i[\psi,\phi]\right)^{\mu\nu}\right)\\
&+\Tr\bigl(i\chi^+_{\mu\nu}[\sigma,\chi^{+\mu\nu}]+
i\chi^-_{\mu\nu}[\sigma,\chi^{-\mu\nu}] \bigr)\biggr).\\\end{aligned}
\end{equation}

Finally,  the useful identity \eqref{eq:ormob}\ enables us to write
% 式 (3.46)
\setcounter{equation}{45}
\begin{equation}\label{eq:nuseful}
\begin{aligned}{{\cmmib I}}_1=-{1\over e^2}\int_Md^4x \sqrt
g&\left(\Tr\left({1\over 2}F_{\mu\nu} F^{\mu\nu}+D_\mu\phi_\nu
D^\mu\phi^\nu +R_{\mu\nu}\phi^\mu\phi^\nu+ {1\over
2}[\phi_\mu,\phi_\nu]^2-(D^*\phi)^2
  \right)\right.\\ &-
2\Tr\left(\chi^+_{\mu\nu}\left(iD\psi+i[\tilde\psi,\phi]\right)^{\mu\nu}
+\chi^-_{\mu\nu}\left(iD\tilde\psi-i[\psi,\phi]\right)^{\mu\nu}\right)\\
&-\Tr\left(i\chi^+_{\mu\nu}[\sigma,\chi^{+\mu\nu}]+
i\chi^-_{\mu\nu}[\sigma,\chi^{-\mu\nu}] \right)\biggr)\\ &
~~~~~~~+{t-t^{-1}\over e^2(t+t^{-1})}\int_M\Tr\, F\wedge F.\\\end{aligned}
\end{equation}

Apart from the topological term, the part of ${{\cmmib
I}}_0+{{\cmmib I}}_1$ that depends only on $A$ and $\phi$ can also
\cite{ref53}\ be written % 式 (3.47)
\setcounter{equation}{46}
\begin{equation}\label{eq:wingo}
{{\cmmib I}}^{(A,\phi)}=-{1\over
e^2}\int d^4x\,\sqrt g\, \Tr\left({1\over 2}
\cF_{\mu\nu}\bar\cF^{\mu\nu}+(D^*\phi)^2\right),
\end{equation}
 as one can show
with some integration by parts similar to what is used in proving
the vanishing theorems.  The analogous terms involving $\sigma$
can be found in \eqref{eq:nobo}\ and are % 式 (3.48)
\setcounter{equation}{47}
\begin{equation}\label{eq:uwingo}
{\cmmib
I}^\sigma={2\over e^2}\int d^4x \sqrt g\,\Tr\left( {1\over
2}[\bar\sigma,\sigma]^2- D_\mu\bar\sigma\,D^\mu\sigma
-[\phi_\mu,\sigma][\phi^\mu,\bar\sigma]\right).
\end{equation}

The key point -- which we already exploited in proving the
vanishing theorems -- is that ${{\cmmib I}}_1$ is independent of
$t$ except for the last term.  But that last term, being a
topological invariant, is automatically annihilated by $\delta_t$,
and indeed by both $\delta_\ell$ and $\delta_\rr$, all by itself.
So without spoiling the topological symmetry, we can add another
term to the action, % 式 (3.49)
\setcounter{equation}{48}
\begin{equation}\label{eq:ucno}
{{\cmmib I}}_2=-\left({t-t^{-1}\over
e^2(t+t^{-1})}-i{\theta\over 8\pi^2}\right)\int_M\Tr\,F\wedge F.
\end{equation}

We have simply chosen the coefficient to cancel the $t$-dependence
in ${{\cmmib I}}_1$, leaving us with the standard $\theta$
parameter multiplying this term in the action.

Finally, then, the overall  action is ${{\cmmib I}}={\cmmib
I}_0+{{\cmmib I}}_1+{{\cmmib I}}_2$. The construction makes it
manifest that ${{\cmmib I}}$ is annihilated by $\delta_t$ for the
specific value of $t$ used in constructing ${{\cmmib I}}_1$.  But
since in fact ${{\cmmib I}}$ is independent of $t$,
it is annihilated by $\delta_t$ for all $t$.

Moreover, ${{\cmmib I}}$ is the action of a topological field
theory. The change in ${{\cmmib I}}$ in a change in the metric of
$M$ is of the form $\delta_t(\dots)$, since ${{\cmmib I}}_0$ and
${{\cmmib I}}_1$ are of this form even before varying the metric,
and ${{\cmmib I}}_2$ does not depend on the metric of $M$ at all.

We have accomplished our goal of constructing an action that has
the full ${\Bbb{ CP}}^1$ family of topological symmetries --
making it clear, among other things, that the partition function
of this theory on a closed four-manifold $M$ without operator
insertions is independent of $t$.  One can readily verify that
(after eliminating $P$ via \eqref{eq:olgog}, whereupon the $(D^*\phi)^2$
terms cancel), the bosonic part of ${{\cmmib I}}$ reduces for
$M=\Bbb{R}^4$ to \eqref{eq:loopy}.  The topological symmetry implies that
the fermionic terms also agree. So this theory is a generalization
of ${\EUN}=4$ super Yang-Mills theory to a curved four-manifold.

\subsection{The Canonical Parameter}

\label{sec:canonpar}

So far, we have obtained a family of topological field theories
that depend on a coupling parameter $\tau=\theta/2\pi+4\pi i/e^2$,
and a twisting parameter $t$.  But these parameters are not really
all independent.  We will now show that the theory really only
depends on a certain combination $\Psi$ of $t$ and $\tau$, which
we will call the canonical parameter.

Let us write the coefficient of the topological term in ${\cmmib
I}_2$, which takes the form $-\left({(t-t^{-1})/
e^2(t+t^{-1})}-i{\theta/8\pi^2}\right)$, as $-\Psi/4\pi i$, where
% 式 (3.50)
\setcounter{equation}{49}
\begin{equation}\label{eq:yorgo}
\Psi = {\tau+\bar\tau\over 2}+{\tau-\bar\tau\over
2}\left({t-t^{-1}\over t+t^{-1}}\right).
\end{equation}
 We will call $\Psi$ the
canonical parameter.

We claim that any observable ${\EUO}(\tau,\bar\tau,t)$ really
depends only on the canonical parameter $\Psi$. Indeed, the action
of the theory is of the form % 式 (3.51)
\setcounter{equation}{50}
\begin{equation}\label{eq:juty}
{{\cmmib I}}=\delta_tV
+{i\Psi\over 4\pi }\int_M\Tr \,F\wedge F.
\end{equation}
 To the extent that the
term $\delta_tV$ is irrelevant, the theory naturally only depends
on the coefficient $\Psi$.

At first sight, there is a fallacy in this argument.  In addition
to whatever $t$-dependence may be explicit in ${{\cmmib I}}$ (we
have arranged so that there is none, though individual terms are
separately $t$-dependent), there might be $t$-dependence coming
from the fact that the definition of $\delta_t$ depends on $t$.
The definition of the theory, after all, depends on the
topological symmetry $\delta_t$ as well as on the action. But in
fact, the $t$-dependence of $\delta_t$ can be eliminated (as long
as $t\not=\pm i$) by redefining the fields. One rescales
$\delta_t$ to $\delta'_t=\delta_t/\sqrt{1+t^2}$, so that the
algebra becomes $(\delta'_t)^2=-i{{\lie}}_\sigma$, independent of
$t$. Then,  with the auxiliary fields $P$ and $H$ included, after
some small redefinitions of the fermions, every multiplet takes
the form $\delta'_t U=V$, $\delta'_tV=-i{{\lie}}_\sigma(U)$, for
some pair of fields $U$, $V$.  (This is a standard type of
multiplet in cohomological field theories, related mathematically
to the Cartan model of equivariant cohomology.)  Once this is
done, the transformation laws of the multiplets are completely
independent of $t$.

\par\medskip{\it Some Special Values}

Some special cases of the relation between $\Psi$, $t$, and $\tau$
are worthy of note. We focus especially on values that have a
simple behavior under $S$-duality and a simple interpretation
after compactification to two dimensions.

If $t=\pm 1$, we have $\Psi={\rm Re}\,\tau$.  Thus, any real value
of $\Psi$ can be achieved at $t=\pm 1$, simply by varying the
$\theta$ angle of the gauge theory, but only real values are
possible. More generally, if $|t|=1$, then $\Psi$ is real.

To make $\Psi$ complex, we must get away from $|t|=1$.  All values
of $\Psi$ in the upper or lower half plane are possible,
respectively, for $t$ real and of modulus greater than or less
than 1.

If $t=\pm i$, then $\Psi=\infty$, independent of $\tau$.  Thus, at
$t=\pm i$, the value of $\tau$ is completely irrelevant.  This has
essentially been demonstrated directly by Marcus \cite{ref53}.

A final observation is that all values of $\Psi$ can be realized
at arbitrarily weak coupling, that is arbitrarily large ${\rm
Im}\,\tau$.  Indeed, inverting the definition of $\Psi$, we have
% 式 (3.52)
\setcounter{equation}{51}
\begin{equation}\label{eq:helmo}
t^2=-{\Psi-\bar\tau\over \Psi-\tau},
\end{equation}
 showing that if
we keep $\Psi$ and ${\rm Re}\,\tau$ fixed and take ${\rm
Im}\,\tau\to\infty$, then $t\to \pm 1$.  The limit ${\rm
Im}\,\tau\to\infty$  is the weak coupling limit where the theory
should be understandable.

\par\medskip{\it Action Of $S$-Duality On The Canonical Parameter}

Now we want to determine how $\Psi$ transforms under $S$-duality.
We claim that a transformation that maps $\tau$ by $\tau\to
(a\tau+b)/(c\tau+d)$ acts likewise on $\Psi$: % 式 (3.53)
\setcounter{equation}{52}
\begin{equation}\label{eq:yrto}
\Psi\to
{a\Psi+b\over c\Psi+d}.
\end{equation}

Clearly, under $T:\tau\to\tau+1$, which leaves $t$ invariant, we
have $\Psi\to \Psi+1$.  We also have to check the behavior under
$S:\tau\to -1/n_{\frak g}\tau$.  This mapping transforms $t$ to
$\pm t\tau/|\tau|$ according to \eqref{eq:duco}. (The sign depends on how
$S$ is lifted to $SL(2,\Bbb{R})$, and does not matter as $\Psi$ is
invariant under $t\to -t$.) So $\Psi$ transforms by
% 式 (3.54)
\setcounter{equation}{53}
\begin{equation}\label{eq:nelmo}
\Psi\to -{1\over 2}\left({1\over n_{\frak
g}\tau}+{1\over n_{\frak g}\bar \tau}\right)-{1\over
2}\left({1\over n_{\frak g}\tau}-{1\over n_{\frak
g}\bar\tau}\right)\left({t^2(\tau/ \bar\tau)-1\over t^2(
\tau/\bar\tau)+1}\right)=-{1\over n_{\frak g} \Psi},
\end{equation}
 where
\eqref{eq:helmo}\ has been used.

\par\medskip{\it Orientation Reversal}

It is a little unusual to find an action of a discrete group such
as $SL(2,\Bbb{Z})$ on a complex parameter $\Psi$ which, as here,
takes values in $\Bbb{CP}^1$, not just in the upper half of the
complex plane.  This being the case, it is conceivable to extend
$SL(2,\Bbb{Z})$ to $GL(2,\Bbb{Z})$ by adding an extra symmetry
that we can take to be % 式 (3.55)
\setcounter{equation}{54}
\begin{equation}\label{eq:ysnon}
{\EUT}=\left(\begin{matrix}-1 & 0 \\
0 & 1\end{matrix}\right).
\end{equation}
  Such an element will act on $\Psi$ by $\Psi\to
-\Psi$, exchanging the upper half plane with the lower half plane.

In fact, such a symmetry exists precisely when the four-manifold
$M$ possesses a symmetry  ${\EUT}:M\to M$  that reverses its
orientation. Such a transformation leaves fixed the gauge coupling
$e^2$ and hence leaves fixed ${\rm Im}\,\tau=4\pi/e^2$. On the
other hand, it reverses the sign of the instanton number and hence
reverses the sign of ${\rm Re}\,\tau$.  Looking back at the
definition of the canonical parameter in \eqref{eq:yorgo}, we see that we
will get $\Psi\to-\Psi$ if ${\EUT}$ acts by % 式 (3.56)
\setcounter{equation}{55}
\begin{equation}\label{eq:zongor}
t\to
{1\over t}.
\end{equation}
 The points $t=\pm 1$ are fixed points of this
transformation.

We will have a symmetry of the topological equations \eqref{eq:loopme}\ if
we take ${\EUT}$ to act on the bosons $(A,\phi)$ by
% 式 (3.57)
\setcounter{equation}{56}
\begin{equation}\label{eq:yoro}
(A,\phi)\to ({\EUT}^*A,-{\EUT}^*\phi).
\end{equation}
 The action on
fermions is % 式 (3.58)
\setcounter{equation}{57}
\begin{equation}\label{eq:zoro}
\begin{aligned}(\psi,\tilde\psi)&\to
({\EUT}^*\tilde\psi,{\EUT}^*\psi)\\ (\chi^+,\chi^-)&\to
({\EUT}^*\chi^-,{\EUT}^*\chi^+)\\ (\eta,\tilde\eta)&\to
({\EUT}^*\tilde\eta,{\EUT}^*\eta).\\\end{aligned}
\end{equation}
  This gives a manifest
symmetry of the action.  The supersymmetry generators transform by
$(u,v)\to(v,u)$ and hence $\delta_t\to t\delta_{t^{-1}}$.

An important application, as we will see, is to the case
$M=\Bbb{R}\times W$, with ${\EUT}$ being a ``time-reversal''
symmetry that acts trivially on $W$ and acts on $\Bbb{R}$ by
multiplication by $-1$.
